\section{实验设置}

\subsection{主线设置}

本文方法实验以 GRU 序列基线为主线。核心设置如下：

\begin{itemize}
  \item 输入模式：feature sequence。
  \item 序列长度：8。
  \item 特征集：manual\_basic，根据任务使用 full 或 compact 子集。
  \item RUL 目标：\code{linear_rul_norm}。
  \item 分类目标：\code{health_state_id} 和 \code{early_fault}。
  \item 训练轮数：200ep。
  \item batch size：64。
  \item 评价指标：RUL 使用 RMSE；分类任务使用 Accuracy、MacroF1 和 WeightedF1，其中 WeightedF1 作为主要指标。
\end{itemize}

所有序列样本均在同一轴承和同一数据子集内部构造，不跨轴承、不跨数据子集。标准化和清洗参数只从训练集估计。

\subsection{对照实验设置}

对照模型包括 MLP、XGBoost 和 RandomForest。它们主要使用窗口级表格特征，不显式建模时间序列。对照实验的作用是判断 manual\_basic 特征本身的建模能力，以及传统表格模型相对于 GRU 序列基线的优势和不足。

\begin{table}[H]
\centering
\caption{对照模型实验的代表性结果}
\label{tab:non-mlp-summary}
\small
\begin{tabularx}{\textwidth}{p{2.5cm}p{2.8cm}p{2.8cm}p{2.4cm}p{2.2cm}Y}
\toprule
数据集 & 任务 & 模型 & 指标 & 测试集结果 & 结论口径 \\
\midrule
XJTU-SY & RUL & RandomForest & RMSE & 0.399658 & 优于默认 MLP，但 RUL 仍有明显分布偏移风险。 \\
XJTU-SY & EarlyFault & XGBoost & WeightedF1 & 0.839369 & 接近 GRU 序列基线，可作为强表格对照。 \\
PHM2012 & RUL & RandomForest & RMSE & 0.337575 & 优于默认 MLP，但弱于 GRU 序列基线。 \\
PHM2012 & EarlyFault & RandomForest & WeightedF1 & 0.613220 & 可作为对照结果，但不是最终主线模型。 \\
\bottomrule
\end{tabularx}
\end{table}


\subsection{演示训练设置}

两个演示视频使用 50ep 训练过程，只用于展示训练过程中的 epoch 变化、train\_loss 变化和日志滚动。它们不是最终指标来源。本文所有主线指标均来自 200ep GRU 序列实验。
